%% kapa: exam=1 round=33 year=2022 period=1 form=A subject=civil_law subject_ko=민법 seq=38 number=38
%% source: Q-Net 감정평가사 기출문제 (cst00309) artlSeq=5209997 file=raw/1008/5209997/2022년도 제33회 감정평가사 1차시험 1교시(민법,경제학,부동산학).pdf
%% conversion: kordoc 4.24.51 -> tools/gen1.py (auto); review=proofread-sonnet; flags=none
\begin{question}{38}
민법 제366조의 법정지상권에 관한 설명으로 옳은 것을 모두 고른 것은? (다툼이 있으면 판례에 따름)
\begin{framedtext}
ㄱ. 미등기건물의 소유를 위해서도 법정지상권이 성립할 수 있다. \par
ㄴ. 당사자 사이에 지료에 관하여 협의한 사실이나 법원에 의하여 지료가 결정된 사실이 없다면, 법정지상권자가 지료를 지급하지 않았다고 하더라도 지료 지급을 지체한 것으로 볼 수 없다. \par
ㄷ. 건물 소유를 위한 법정지상권을 취득한 사람으로부터 경매에 의해 건물소유권을 이전받은 매수인은 특별한 사정이 없는 한 건물의 매수취득과 함께 위 지상권도 당연히 취득한다.
\end{framedtext}
\begin{choices}
  \item ㄱ
  \item ㄴ
  \item ㄱ, ㄷ
  \item ㄴ, ㄷ
  \item ㄱ, ㄴ, ㄷ
\end{choices}
\end{question}
\answer{5}
